\documentclass[11pt]{article}
\usepackage[margin=1in]{geometry}
\usepackage{amsmath,amssymb,amsthm}
\usepackage[hidelinks]{hyperref}

\newcommand{\T}{\mathbb T}

\title{Literature Status: Hypercontractive Homogeneous Cotype}
\author{Run \texttt{fa\_banach\_001}, agent \texttt{agent\_lane\_06}}
\date{2026-06-28}

\begin{document}
\maketitle

\section*{Source Question}

The source paper is Daniel Carando, Andreas Defant, and Pablo Sevilla-Peris,
\emph{Some polynomial versions of cotype and applications},
arXiv:1503.00850.

In Section 1, after defining hypercontractive homogeneous cotype, the authors ask whether this property is equivalent to ordinary cotype. In their notation, a Banach space $X$ has hypercontractive homogeneous cotype $q$ if there is $C>0$ such that for every $m$ and every finitely supported family $(x_\alpha)_{|\alpha|=m}\subset X$,
\[
\left(\sum_{|\alpha|=m}\|x_\alpha\|^q\right)^{1/q}
\le C^m
\left(
\int_{\T^\infty}
\left\|
\sum_{|\alpha|=m} x_\alpha z^\alpha
\right\|^2\,dz
\right)^{1/2}.
\]
They conjecture that, for each $q\ge 2$, a Banach space has this property if and only if it has ordinary Rademacher cotype $q$.

\section*{Answering Literature}

The later paper is Daniel Carando, Felipe Marceca, and Pablo Sevilla-Peris,
\emph{Hausdorff-Young type inequalities for vector-valued Dirichlet series},
arXiv:1904.00041.

Its introduction explicitly states that the paper proves Rademacher cotype is equivalent to the hypercontractive homogeneous cotype defined in arXiv:1503.00850. The precise theorem is Theorem \texttt{cotimplicapolcot} in Section 3. It proves that, for $2\le q<\infty$, the following are equivalent:
\begin{itemize}
\item $X$ has cotype $q$;
\item there is $C>0$ such that every $X$-valued polynomial of degree at most $m$ satisfies
\[
\left(\sum_{|\alpha|\le m}\|x_\alpha\|^q\right)^{1/q}
\le C^m
\left(
\int_{\T^n}
\left\|
\sum_{|\alpha|\le m}x_\alpha z^\alpha
\right\|^q dz
\right)^{1/q}.
\]
\end{itemize}
Remark \texttt{expo} then records that the polynomial $L_q$ norm can be replaced by any $L_r$ norm, with constants still growing only exponentially in $m$. Restricting this theorem to homogeneous polynomials and using $r=2$ gives the definition of hypercontractive homogeneous cotype in the source paper. The converse direction is immediate already from the case $m=1$.

\section*{Status}

The conjecture from arXiv:1503.00850 is therefore already answered affirmatively by arXiv:1904.00041. This packet is a literature-status packet, not a new proof produced by the run.

\begin{thebibliography}{9}
\bibitem{CDS2015}
D. Carando, A. Defant, and P. Sevilla-Peris,
\emph{Some polynomial versions of cotype and applications},
arXiv:1503.00850, 2015.

\bibitem{CMSP2019}
D. Carando, F. Marceca, and P. Sevilla-Peris,
\emph{Hausdorff-Young type inequalities for vector-valued Dirichlet series},
arXiv:1904.00041, 2019.
\end{thebibliography}

\end{document}
